% GENERATED by scripts/manual_tables.py from the adapter registry. Do not edit:
% rerun the script. Hand-written text: blocks/detection/spike_v1/
\ubblock{detection.spike_v1}%
  {Spike detection (v1)}% what the app calls it
  {Spike detection (Adamatzky 2023)}% registry display name
  {detect}% insert-modal tab
  {Signal}% input
  {SpanSet}% output
  {Adapters/detection\_spike\_v1.py}% code
  {detection/spike_v1}% hand-written folder under blocks/
\ubblockdescription{detection.spike_v1}{Compares every sample with the average of its neighbours and calls it a spike when it stands out from that local average by more than delta\_mv (in mV) - tune delta\_mv first; samples closer than d are one spike. Raise delta\_mv if it fires on noise, lower it if it misses spikes you can see; then width (how local the average is; it is blind to changes slower than 4w+1 samples, so use Drop detection for slow drops), then d. Each span is one spike, marked inside it.}
\ubblockflags{detection.spike_v1}{%
  \ubflag{units}{}%
  \ubflag{derive}{}%
}
\ubblocksideinputs{detection.spike_v1}{%
}
\ubblockparams{detection.spike_v1}{%
  \ubparam{delta\_\allowbreak{}mv}{1.0}{number, at least 0}{Tune this first. How far (mV) a sample must stand from its local average to be a spike; converted by the recording's declared unit. 1 mV is a placeholder that holds at M2's noise (the paper's 0.01 is in its logger's mV and fires on M2's noise); raise it on noisier channels - the readout gives it as a multiple of this span's own wander}%
  \ubparam{width}{20}{whole number, at least 1}{Neighbourhood half-width w (samples): the local average spans 4w+1 samples; changes slower than that are invisible to this test (paper: 20)}%
  \ubparam{d}{30}{whole number, at least 1}{Refractory gap (samples): candidates closer than this are one spike, at its most extreme sample (paper: 30)}%
  \ubparam{polarity}{\texttt{both}}{one of \texttt{both}, \texttt{up}, \texttt{down}}{Which direction counts: both, up (rises) or down (dips)}%
  \ubparam{span\_\allowbreak{}s}{0.0}{number, at least 0}{Width of the span drawn round each spike, seconds; 0 = the refractory gap d}%
}
